\documentclass[10pt,a4paper]{amsart}
\usepackage[margin=19mm]{geometry}
\usepackage{amsmath,amssymb,mathtools}
\usepackage[scr=rsfs,cal=euler]{mathalfa}
\usepackage{xfrac}
\usepackage[unicode,hidelinks,bookmarks=false]{hyperref}
\newcommand{\ie}{i.e.\ }
\newcommand{\cf}{cf.\ }
\newcommand{\resp}{respectively\ }
\newcommand{\Subsection}{\subsection}
\providecommand{\nobreakdash}{\nobreak\hbox{-}}
\setlength{\emergencystretch}{2em}
\multlinegap0pt
\usepackage{amssymb}
\usepackage{bm}
\newtheorem{definition}{Definition}
\newtheorem{lemma}{Lemma}
\newcommand{\W}{\mathcal{W}}
\title{Cliques in graphs constructed from Strongly Orthogonal Subsets in exceptional root systems}
\author{Patrick J. Browne, P\'adraig \'O Cath\'ain}
\date{Source: arXiv:2604.02983v1 (CC BY 4.0)}
\begin{document}
\maketitle

\noindent\emph{Source and licence.} Cliques in graphs constructed from Strongly Orthogonal Subsets in exceptional root systems. Version arXiv:2604.02983v1 (CC BY 4.0), distributed by arXiv under the Creative Commons Attribution 4.0 International licence, http://creativecommons.org/licenses/by/4.0/, as recorded in the arXiv licence metadata for this exact version and shown on the abstract page https://arxiv.org/abs/2604.02983v1. Full licence notice: \texttt{licenses/arxiv-2604.02983-CC-BY-4.0.txt}.\par\smallskip

\noindent\textit{Editorial scope.} Counted unit: the article's lemma lem:mod8 (source0.tex lines 180--182). The article's complete original proof of it is delivered with it (source0.tex lines 184--197, one \texttt{proof} environment of 14 lines). No mathematics is written, repaired or re-derived: the statement and the proof are the article's own lines, in the article's own notation, and no retained line is substituted or re-flowed.  Retained uncounted as the source-local context the proof needs: lines 57--67.  Nothing was omitted from any retained span, so the package carries no omission record.  No retained line contains a citation, so the wrapper carries no bibliography and no bibliography entry.  No retained line contains a figure environment, an image-inclusion command or drawing code, so no image is delivered and the package declares no asset dependency.  Editorial formatting only, in addition: an AMS \texttt{amsart} preamble that restates the article's own declarations and macros used by the retained lines, namely the environments \texttt{definition}, \texttt{lemma} on separate counters, together with the standard packages those lines need.  The retained environments print in this document as Definition 1, Lemma 1 in order of appearance, whereas the article prints the same items with the numbering of its own full document.  The complete native source of the article as distributed in the arXiv e-print for this exact version is bundled unchanged as \texttt{sources/197772/source0.tex}.
\begin{definition}\label{RootDefn}
Let $V$ be a finite-dimensional Euclidean space with inner product $\langle \cdot, \cdot \rangle$.
A \emph{root system} in $V$ is a finite set $R \subset V$ satisfying:
\begin{enumerate}
    \item $R$ spans $V$ and $0 \notin R$.
    \item If $\alpha \in R$, then $-\alpha \in R$ and no other scalar multiple of $\alpha$ lies in $R$.
    \item For all $\alpha, \beta \in R$, the reflection $s_\alpha(\beta) = \beta - 2\frac{\langle \beta, \alpha \rangle}{\langle \alpha, \alpha \rangle}\alpha$ lies in $R$.
    \item For all $\alpha, \beta \in R$, we have $2\frac{\langle \beta, \alpha \rangle}{\langle \alpha, \alpha \rangle} \in \mathbb{Z}$.
\end{enumerate}
The \emph{rank} of $R$ is $\dim V$. The \emph{Weyl group}, $\W(R)$, is the subgroup of $\mathrm{GL}(V)$ generated by the reflections $\{s_\alpha : \alpha \in R\}$. It is a finite group that permutes $R$ and acts linearly on $V$.
\end{definition}

\begin{lemma}\label{lem:mod8}
Let $R$ be a simply-laced root system of rank $\ell$, and suppose $k = \ell$. Then for any two vertices $v, w \in \mathcal{V}(R, k)$, we have $\|v - w\|^2 \equiv 0 \pmod{8}$.
\end{lemma}

\begin{proof}
Write $v = \sum_{i=1}^{\ell} \alpha_i$ and $w = \sum_{j=1}^{\ell} \beta_j$, where $\{\alpha_i\}$ and $\{\beta_j\}$ are $\ell$-element strongly orthogonal subsets. Since the $\alpha_i$ are pairwise orthogonal vectors of norm $2$ in a space of dimension $\ell$, they satisfy $\sum_{i=1}^{\ell} \langle \alpha_{i}, v\rangle v = 2v$ for any vector $v$; and similarly for $\{\beta_j\}$. By Axiom 4 of Definition \ref{RootDefn}, the entries of the matrix  $C$ given by $c_{ij} = \langle \alpha_i, \beta_j \rangle$ are integers. A computation shows that $C^\top C = 4 I_\ell$, and that $CC^\top = 4I_\ell$. 
So the rows of $C$ have squared entries summing to $4$: each row contains either a single non-zero entry $\pm 2$, or four non-zero entries $\pm 1$; in all cases the row sum is even.

Let $\mathbf{1} = (1, \ldots, 1)^\top$; and observe that the vector $s = \frac{1}{2}C\mathbf{1}$ has integer entries. Now
\[
\langle v, w \rangle = \mathbf{1}^\top C \mathbf{1} = 2 \sum_{i} s_{i},
\]
From the definitions of $s$ and $C$, we deduce that $\|s\|^2 = \tfrac{1}{4}\mathbf{1}^\top C^\top C \mathbf{1} = \tfrac{1}{4} \cdot 4\ell = \ell$. 
Since $n^2 \equiv n \pmod{2}$ for any integer $n$, we obtain $\sum_i s_i \equiv \sum_i s_i^2 = \ell \pmod{2}$. Therefore $\langle v, w \rangle = 2\sum_i s_i \equiv 2\ell \pmod{4}$, and
\[
\|v - w\|^2 = \|v\|^2 + \|w\|^2 - 2\langle v, w \rangle = 4\ell - 2\langle v, w \rangle \equiv 4\ell - 4\ell = 0 \pmod{8}. \qedhere
\]
\end{proof}

\end{document}
